\section{讨论}
\label{sec:discussion}

\paragraph{空间感知为何重要。}全局平均池化是视觉运动策略中最常见的特征聚合方式，但它将$512 \times H' \times W'$的丰富特征图压缩为单个512维向量。对于粗粒度任务（如抓取大型物体），这种信息丢失可能无关紧要。但对于精密任务（如焊缝跟踪、插入、缝合），空间信息至关重要——策略需要知道工具尖端在哪里、焊缝在哪里、两者之间的相对位置关系如何。\fasa{}通过自监督变化检测恢复了这些空间信息，且仅增加极少的参数开销。

\paragraph{自监督 vs. 基于蒸馏的空间注意力。}视觉基础模型蒸馏（如从Grounding DINO或CLIP）可以提供空间注意力监督，但存在重要局限：（1）需要访问特定的教师模型和离线特征提取流水线；（2）蒸馏的注意力是\emph{静态的}——反映教师的通用目标检测能力，而非特定任务的动态交互模式；（3）不能跨领域泛化，除非重新运行教师模型。\fasa{}通过从任务本身的时间结构中推导空间注意力来解决所有三个局限。基于变化的注意力自然适应当前相关的内容：接近时，它突出显示在静态场景中移动的机器人；跟踪时，它聚焦于微妙的工具-工件接触区域。在我们的框架中，可选的预蒸馏增强了主干网络的特征质量（有利于\fasa{}的变化检测），而\fasa{}在池化层面提供动态的、任务自适应的空间注意力。

\paragraph{焊缝跟踪之外的适用性。}虽然我们在焊缝跟踪任务上评估\fasa{}，但其核心机制——帧间特征变化检测并据此分配空间注意力——是完全任务无关的。任何需要细粒度空间理解的精密任务——如插入、缝合或接触密集型操作——都可以从\fasa{}中受益。\fasa{}的自监督特性意味着迁移到新领域时不需要任务特定的标注或重新训练教师模型。

\paragraph{计算开销。}表~\ref{tab:compute}总结了参数和推理时间开销。\fasa{}模块每个相机仅增加约67k参数（两个相机共约134k），就提供了$2.07\times$更丰富的空间特征（1060D vs. 512D），是一个极其高效的升级。相比基础VITA的约31.9M参数，\fasa{}仅增加约0.4\%的参数量。

\begin{table}[t]
\centering
\caption{\fasa{}相对于基础VITA策略的计算开销。}
\label{tab:compute}
\begin{tabular}{lcc}
\toprule
\textbf{组件} & \textbf{参数量} & \textbf{推理时间 (ms)} \\
\midrule
基础VITA（全局平均池化） & $\sim$31.9M & -- \\
\midrule
\multicolumn{3}{l}{\textit{\fasa{}模块}} \\
\quad 变化精炼网络（每相机） & +$\sim$34k & +$<$0.3 \\
\quad 坐标投影（每相机） & +$\sim$33k & +$<$0.1 \\
\quad obs\_encoder 维度增大 & +$\sim$0.5M & +$<$0.1 \\
\midrule
\textbf{完整\method{}} & $\sim$32.4M (+1.6\%) & +$<$0.5 \\
\bottomrule
\end{tabular}
\end{table}

\paragraph{局限性。}以下局限性需要说明：
\begin{itemize}[leftmargin=*,itemsep=2pt]
    \item \textbf{\fasa{}冷启动。}\fasa{}需要$\text{obs\_horizon} \geq 2$来计算帧间变化。在每个episode的第一步（无前一帧），它降级为全局池化加空间softmax。对于第一个动作至关重要的任务，这可能影响初始性能。
    \item \textbf{依赖视觉变化。}\fasa{}的注意力基于帧间特征变化。在极慢运动或静止阶段，变化信号可能过弱，导致注意力退化为近似均匀分布。不过在实践中，50~Hz采样下即使缓慢运动也产生可检测的特征变化。
    \item \textbf{单任务评估。}我们在单个任务领域上进行评估。虽然\fasa{}设计为通用的即插即用模块，但在多样化操作任务上的更广泛评估将加强我们的主张。
\end{itemize}

\paragraph{未来方向。}有前景的扩展包括：（1）在真实机器人硬件上部署\method{}以验证空间感知的仿真到真实迁移；（2）探索在多个ResNet层输出上操作的多尺度\fasa{}以实现层次化空间注意力；（3）将\fasa{}的变化注意力图与好奇心驱动探索结合用于强化学习；（4）将\fasa{}扩展到多帧时间对比（超越两个连续帧）以实现更丰富的运动模式检测；（5）探索将时间感知通过非视觉通道（如关节速度、动作历史）安全地注入策略，避免干扰视觉特征学习。
